%% kapa: exam=1 round=37 year=2026 period=1 form=A subject=civil_law subject_ko=민법 seq=2 number=2
%% source: Q-Net 감정평가사 기출문제 (cst00309) artlSeq=5258924 file=raw/1008/5258924/2026년 제37회 감정평가사 1차 1교시.pdf
%% conversion: kordoc 4.24.51 -> tools/gen1.py (auto); review=proofread-sonnet; flags=none
\begin{question}{2}
부재자의 재산관리에 관한 설명으로 옳지 않은 것은? (다툼이 있으면 판례에 따름)
\begin{choices}
  \item 법원이 선임한 재산관리인은 관리할 재산목록을 작성하여야 한다.
  \item 법원은 그 선임한 재산관리인에 대해 부재자의 재산으로 상당한 보수를 지급할 수 있다.
  \item 재산관리인의 부재자 재산의 처분행위에 대한 법원의 허가는 과거의 처분행위를 추인하는 방법으로는 할 수 없다.
  \item 부재자의 생사가 분명하지 아니한 경우, 부재자가 정한 재산관리인이 권한을 넘은 행위를 할 때에는 법원의 허가를 받아야 한다.
  \item 법원이 선임한 재산관리인이 그 선임 취소 전에 한 권한 내의 행위는 부재자에 대한 실종선고기간 만료 후에 이루어진 경우에도 유효하다.
\end{choices}
\end{question}
\answer{3}
